% ============================================================
% SKELETON -- Part II, Section 7: The Ideal Gas
% ============================================================

\section{The Ideal Gas}

\paragraph{Scope of this section.}
The \textcolor{blue}{ideal gas} is the first explicit model of the notes.
The model consists of point particles without interactions. The canonical
ensemble of Section~3 gives every property of the model, so the ideal gas
is a result of the assumptions rather than an empirical law.

\subsection{The Classical Partition Function}

\paragraph{From a sum over microstates to a phase-space integral.}
Classical particles have continuous positions and momenta, so the sum over
microstates becomes an integral over phase space. Two assumptions enter.
Each microstate occupies a phase-space volume $h^{3N}$, and the factor
$1/N!$ removes the permutations of identical particles.
% FROM: 1/N! and extensivity (Section 5, Gibbs paradox).
% TODO: state the h^{3N} cell as an assumption, justified by quantum
%       statistics in Part V.
\begin{equation}
    \boxed{Z=\frac{1}{N!\,h^{3N}}\int d^{3N}q\,d^{3N}p\;e^{-\beta\mathcal{H}(q,p)}.}
\end{equation}

\paragraph{Partition function of non-interacting particles.}
The Hamiltonian contains the kinetic energy only. The momentum integrals
are Gaussian, and the position integrals give the volume.
\textcolor{red}{Here, $\lambda\equiv h/\sqrt{2\pi mk_BT}$ denotes the
thermal wavelength.}
\begin{equation}
    \boxed{Z=\frac{V^N}{N!\,\lambda^{3N}}.}
\end{equation}

\subsection{Thermodynamics of the Ideal Gas}

\paragraph{Equation of state, energy and heat capacity.}
The free energy $F=-k_BT\ln Z$ gives the pressure, and
$U=-\partial\ln Z/\partial\beta$ gives the energy.
% FROM: F = -k_B T ln Z and P = -dF/dV (Section 4).
\begin{equation}
    \boxed{PV=Nk_BT,\qquad U=\tfrac{3}{2}Nk_BT,\qquad C_V=\tfrac{3}{2}Nk_B.}
\end{equation}

\paragraph{Entropy and chemical potential.}
The entropy of the ideal gas is the \textcolor{blue}{Sackur--Tetrode
equation}. The entropy is extensive only with the factor $1/N!$.
\begin{equation}
    \boxed{S=Nk_B\left[\ln\frac{V}{N\lambda^3}+\frac{5}{2}\right],\qquad
    \mu=k_BT\ln\!\left(n\lambda^3\right).}
\end{equation}

\paragraph{Absolute activity.}
The grand canonical weight $e^{\beta\mu N}$ defines the
\textcolor{blue}{absolute activity} $z\equiv e^{\beta\mu}$. For the ideal
gas, the absolute activity is the number of particles in a volume
$\lambda^3$.
% FROM: Xi(beta, mu) (Section 5).
% TODO: Xi = exp(zV/lambda^3); N = zV/lambda^3.
\begin{equation}
    \boxed{z=e^{\beta\mu}=n\lambda^3.}
\end{equation}

\paragraph{Chemical potential and pressure.}
At fixed temperature, the chemical potential of the ideal gas depends on
the pressure through a logarithm. This form is the reference for the
fugacity of a real gas in Section~8.
\textcolor{red}{Here, $P^\circ$ denotes a reference pressure, usually
1~bar, and $\mu^\circ(T)$ the chemical potential at $P^\circ$.}
\begin{equation}
    \boxed{\mu(T,P)=\mu^\circ(T)+k_BT\ln\frac{P}{P^\circ}.}
\end{equation}

\paragraph{Response functions of the ideal gas.}
The general results of Section~6 take simple forms. The particle-number
fluctuation equals the mean, as for a Poisson distribution.
% TODO: kappa_S = 3/(5P); C_P = 5/2 N k_B; gamma = C_P/C_V = 5/3.
\begin{equation}
    \boxed{\alpha_P=\frac{1}{T},\qquad\kappa_T=\frac{1}{P},\qquad
    C_P-C_V=Nk_B,\qquad\frac{\mathrm{Var}(N)}{N}=1.}
\end{equation}

\paragraph{Mixtures of ideal gases.}
In a mixture of ideal gases, each species behaves as if alone in the
volume. The \textcolor{blue}{partial pressure} $P_i=x_iP$ then replaces $P$
in the chemical potential of species $i$.
\textcolor{red}{Here, $x_i\equiv N_i/N$ denotes the mole fraction.}
% FROM: mixing entropy (Section 5).
% TODO: Dalton's law; mixing entropy -N k_B sum x_i ln x_i.
\begin{equation}
    \boxed{P=\sum_iP_i,\qquad
    \mu_i=\mu_i^\circ(T)+k_BT\ln\frac{x_iP}{P^\circ}.}
\end{equation}

\paragraph{Validity of the ideal gas.}
The ideal gas requires two conditions. The gas must be dilute, so that the
interactions are negligible, and the thermal wavelength must be small
compared with the particle spacing, $n\lambda^3\ll1$, so that the classical
limit holds.
